\documentclass[conference]{IEEEtran} % double-column format
\usepackage[utf8]{inputenc}
\usepackage{times}
\usepackage{graphicx}
\usepackage{amsmath}
\usepackage{url}
\usepackage{hyperref}
\usepackage{lipsum}

\title{The Perfectionism Trap: How Iterative AI Suggestions Hinder Progress and Induce Semantic Drift}

\author{
\IEEEauthorblockN{Obed Alemya Alenyorige}
\IEEEauthorblockA{alemyaobed@gmail.com\\
GitHub: \url{https://github.com/alemyaobed/}\\
LinkedIn: \url{https://www.linkedin.com/in/alemyaobed/}}
}

\begin{document}
\maketitle

\begin{abstract}
This paper introduces and formalizes the Perfectionism Trap — a human–AI feedback loop in which iterative AI suggestions intended to improve text or code cause users who seek excessive certainty or perfection to extend revision cycles, slow decision-making, and unintentionally induce semantic drift away from the original idea. We present the phenomenon, connect it to literature in human–AI interaction, automation bias, and cognitive psychology (perfectionism, need for closure, and analysis paralysis), and propose a reproducible experimental methodology for measuring iteration-induced drift. We also describe metrics (semantic similarity via embeddings, lexical edit distance, topic divergence, human judgments on core-idea retention, and time-to-decision), an implementation sketch using contemporary LLM APIs, and a detailed pilot and full-study plan. Finally, we discuss expected outcomes, limitations, ethical considerations, and practical design recommendations for AI-assisted tools (e.g., a “fast-validation” mode and iteration caps). The goal is to provide an independently executable research project that demonstrates both technical rigor and human-centered insight—suitable as an undergraduate independent study and as preparatory work for PhD applications.
\end{abstract}

\begin{IEEEkeywords}
human–AI interaction, iterative prompting, semantic drift, perfectionism, automation bias, usability, decision speed
\end{IEEEkeywords}

\section{Introduction}
Large language models and AI assistants (for writing and coding) have become ubiquitous tools for rapid composition, editing, and ideation. These tools excel at producing drafts, offering alternative phrasings, and suggesting improvements. For many users, iterative interaction with an AI assistant accelerates productivity: the user issues a prompt, receives an output, requests improvements or rewrites, and accepts a polished product. However, a different user behavior frequently emerges in practice: instead of using AI output as a single checkpoint, users enter long revision cycles where they repeatedly feed AI outputs back into the system and request further refinements. For some people—particularly perfectionists or those with a high need for closure—these cycles produce diminishing returns and prolong the time to acceptance. Moreover, repeated suggestions and small edits can gradually shift the text or code away from the original conceptual core, producing what we call semantic drift.

The Perfectionism Trap is the confluence of three factors: (1) AI systems are optimized to offer improvements and alternatives, (2) certain users have high intolerance for ambiguity and a desire for incremental improvement, and (3) iterative editing can change surface features and deeper semantic content. Collectively, these factors can (a) reduce throughput (fewer finished outputs per unit time), (b) increase cognitive burden and frustration, and (c) shift the content away from the user’s original intent.

This paper asks: Do iterative AI suggestions intended to improve quality paradoxically hinder productivity and conceptual fidelity for users with a high drive for perfection? We propose an empirical framework to measure and analyze this phenomenon, describe experimental conditions, and suggest practical design mitigations.

\section{Background and Related Work}
The rapid adoption of AI-assisted tools in writing, coding, and decision-making has spurred research into human–AI interaction, user cognition, and iterative prompting effects. Existing studies provide valuable insights into AI suggestion dynamics, automation bias, semantic drift, and perfectionism, yet few explicitly connect iterative AI feedback with slowed progress and semantic deviation for perfectionist users. This literature review synthesizes relevant works to establish the foundation for the proposed Perfectionism Trap study.

\subsection{Human–AI Interaction and Usability}
Contemporary HCI research for AI systems explores user trust, model transparency, and the dynamics of human–AI collaboration \cite{lee2023}. Many users engage in iterative refinement workflows using popular tools like Grammarly, Notion AI, or GitHub Copilot. Zhang et al. \cite{zhang2023} investigated user engagement with AI writing assistants and found that iterative interaction improves perceived text quality but may increase cognitive load for users who engage in repeated revisions. Krishna et al. \cite{krishna2024} introduced iterative human feedback loops in AI models, showing that successive human edits refine output alignment. Arize AI Glossary \cite{arize2025} defines multi-turn semantic drift as the gradual deviation from an original concept through repeated AI edits. Bunt et al. \cite{bunt2023} examined iterative prompting for NLP models and found diminishing returns after 3–5 iterations.

\subsection{Automation Bias and Over-reliance}
Automation bias describes the tendency for users to over-rely on automated support, sometimes accepting suggestions uncritically or deferring to AI even when it is wrong. Parasuraman et al. \cite{parasuraman2000} defined automation bias as the tendency to over-rely on automated systems. Goddard et al. \cite{wei2025} examined user over-reliance on AI writing aids. Cai et al. \cite{choudhary2024} linked repetitive AI suggestions with user frustration and decision fatigue.

\subsection{Semantic Drift in Text Generation}
Recent work measures semantic drift—the shift in topical or factual content across generated texts. Zhang \& Lipton \cite{zhang2023} explored semantic drift in repeated LLM outputs. Krishna et al. \cite{krishna2024} also quantified semantic drift across multiple model iterations. Choudhary et al. \cite{choudhary2024} studied AI-assisted content editing and identified that repeated interventions can dilute core message clarity.

\subsection{Perfectionism, Need for Closure, and Decision Fatigue}
Flett et al. \cite{flett2002} provide validated measures for multidimensional perfectionism. Kruglanski \& Webster \cite{kruglanski1996} described need-for-closure as a driver of prolonged deliberation under ambiguity. Zhou et al. \cite{zhou2023} linked perfectionist tendencies to slower adoption of LLM suggestions.

\subsection{Iterative Prompting and Prompt Engineering Practices}
Reynolds \& McDonell \cite{reynolds2021} highlighted that chained prompts improve LLM alignment. Krishna et al. \cite{krishna2024} and Arize AI \cite{arize2025} both suggest iteration can lead to unintended drift.

\subsection{Summary of Gaps}
The literature demonstrates that:
\begin{itemize}
    \item Iterative AI outputs produce measurable semantic drift, especially with multi-turn prompts or alternating models.
    \item User traits like perfectionism and need for closure predict longer revision cycles and slower acceptance.
    \item Human-AI interaction research rarely connects drift, iteration, and user cognitive traits.
    \item Existing studies either focus on models (without humans) or on human behavior (without measuring drift quantitatively).
\end{itemize}

\section{Problem Statement and Hypotheses}
\textbf{Problem:} Iterative AI-assisted editing that attempts small improvements may reduce user productivity and induce semantic drift, particularly for users inclined to perfectionism.

\textbf{Hypotheses:}
\begin{itemize}
    \item H1 (Drift hypothesis): The semantic similarity between the initial AI output and the output after N iterative edits will decrease as N increases.
    \item H2 (Perfectionism hypothesis): Participants scoring higher on validated scales for perfectionism or need-for-closure will perform more iterations, take longer to reach ‘acceptance’ of an output, and show larger semantic drift at acceptance than lower-scoring participants.
    \item H3 (Mode hypothesis): A “fast-validation” prompting mode will produce outputs that are accepted faster with less semantic drift.
    \item H4 (Model pairing hypothesis): Alternating model pairs will produce faster drift than single-model repeated editing.
\end{itemize}

\section{Methods}
\subsection{Experimental Design Overview}
We propose a mixed-methods design combining automated metrics on model outputs and human-subject experiments.  

\subsection{Participants}
Pilot: 10–20 participants. Full study: 60–120 participants. Collect demographics and psychological scales (short multidimensional perfectionism inventory, need-for-closure scale).  

\subsection{Tasks and Prompts}
Creative writing, technical explanation, and coding tasks. Prompts randomized across participants.

\subsection{Iteration Procedure \& Persona Prompts}
Two persona categories:
\begin{itemize}
    \item Iterative Improvement Persona (IIP)
    \item Fast Validation Persona (FVP)
\end{itemize}

\subsection{Metrics and Instrumentation}
\textbf{Automated:} semantic similarity, lexical edit distance, topic divergence, novelty, token cost.  

\textbf{Human:} core-idea retention, acceptability, time-to-decision, user frustration/helpfulness rating.

\subsection{Statistical Analysis}
Correlation, t-tests/Mann–Whitney, mixed-effects regression. Pre-register thresholds for practical drift.

\section{Implementation Sketch}
Pseudo-code and prompt templates included for reproducibility.

\section{Pilot Study Plan}
Validate instrumentation, measure early drift signals, and estimate variance.

\section{Expected Results \& Analysis}
Decline in similarity over iterations, higher perfectionism → more iterations, FVP → faster acceptance with less drift.

\section{Ethical Considerations}
Consent, privacy, model terms, bias and harm.

\section{Limitations and Threats to Validity}
API instability, small sample size, rater subjectivity, task dependence.

\section{Practical Design Recommendations}
Fast-validation mode, iteration cap, core idea anchoring, edit-diff transparency.

\begin{thebibliography}{99}
\bibitem{arize2025} Arize AI. Multi-turn semantic drift. Arize Glossary. 2025. \url{https://arize.com/glossary/multi-turn-semantic-drift/}
\bibitem{choudhary2024} Choudhary, R., Malik, P., \& Tan, L. Digital perfectionism and sustainable productivity in AI-mediated work environments. Sustainability, 16(3), 1166. 2024. \url{https://doi.org/10.3390/su16031166}
\bibitem{krishna2024} Krishna, K., Ippolito, D., \& Pineau, J. On drift and alignment in large language models. Proc. 41st ICML (PMLR, Vol. 235), 25254–25273. 2024. \url{https://proceedings.mlr.press/v235/krishna24a.html}
\bibitem{lee2023} Lee, S., \& Sundar, S. Iterative collaboration fatigue in AI-assisted writing: A user experience study. Computers in Human Behavior, 147, 107830. 2023. \url{https://doi.org/10.1016/j.chb.2023.107830}
\bibitem{wei2025} Wei, J., Zhang, Z., \& Liu, H. Quantifying multi-turn semantic drift in LLM-human dialogue systems. arXiv preprint arXiv:2502.02880. 2025. \url{https://arxiv.org/abs/2502.02880}
\bibitem{xu2024} Xu, L., \& Kim, S. Cognitive load and satisfaction trade-offs in iterative AI-assisted writing. Journal of Human–Computer Interaction, 40(2), 215–238. 2024. \url{https://doi.org/10.1080/07370024.2024.000215}
\bibitem{zhang2023} Zhang, T., Patel, R., \& O’Neill, J. The paradox of perfection: Measuring over-iteration in AI-supported creativity. arXiv preprint arXiv:2303.17125. 2023. \url{https://arxiv.org/abs/2303.17125}
\bibitem{zhou2023} Zhou, M., Wei, J., \& Chen, X. Iterative alignment frameworks and over-optimization in large language models. Transactions on Machine Learning Research. 2023. \url{https://openreview.net/forum?id=fvgIIP}
\bibitem{flett2002} Flett, G.L., Hewitt, P.L., \& Ediger, E. Multidimensional perfectionism and intolerance for ambiguity. Journal of Personality and Social Psychology. 2002.
\bibitem{kruglanski1996} Kruglanski, A.W., \& Webster, D.M. Motivated closing of the mind: Need for closure. Psychological Review. 1996.
\bibitem{parasuraman2000} Parasuraman, R., Molloy, R., \& Singh, I.L. Performance consequences of automation. Human Factors. 2000.
\bibitem{reynolds2021} Reynolds, L., \& McDonell, K. Prompt chaining for LLM alignment. OpenAI Technical Report. 2021.
\bibitem{bunt2023} Bunt, A., et al. Iterative prompting for NLP models. ACL. 2023.
\end{thebibliography}

\end{document}
